%% kapa: exam=1 round=29 year=2018 period=1 form=A subject=economics subject_ko=경제학원론 seq=14 number=54
%% source: Q-Net 감정평가사 기출문제 (cst00309) artlSeq=5202808 file=extracted/1008/5202808/2018년 제29회 감정평가사 1차시험 기출문제/2018년 제29회 감정평가사 1차시험 기출문제/1교시 A형.pdf
%% conversion: kordoc 4.24.51 -> tools/gen1.py (auto); review=proofread-sonnet; flags=none
\begin{question}{54}
기업 A의 생산함수는 $Q=\min\{L,K\}$이다. 이에 관한 설명으로 옳은 것을 모두 고른 것은? (단, $Q$는 산출량, $w$는 노동 $L$의 가격, $r$은 자본 $K$의 가격이다.)
\begin{framedtext}
ㄱ. 생산요소 $L$과 $K$의 대체탄력성은 0이다. \par
ㄴ. 생산함수는 1차 동차함수이다. \par
ㄷ. 비용함수는 $C(w,r,Q)=Q^{w+r}$로 표시된다. \par
\end{framedtext}
\begin{choices}
  \item ㄱ
  \item ㄴ
  \item ㄱ, ㄴ
  \item ㄴ, ㄷ
  \item ㄱ, ㄴ, ㄷ
\end{choices}
\end{question}
\answer{3}
